\documentclass[11pt,a4paper]{altacv}

\geometry{left=1.5cm,right=1.5cm,top=1.cm,bottom=1.2cm,columnsep=1.5cm}

\usepackage[utf8]{inputenc}
\usepackage[T1]{fontenc}
\usepackage[french]{babel}
\usepackage{paracol}
\usepackage{xcolor}
\usepackage{hyperref}
\usepackage{fontawesome5}
\usepackage{setspace}

\definecolor{SlateGrey}{HTML}{2E2E2E}
\definecolor{LightGrey}{HTML}{666666}
\definecolor{DarkPastelRed}{HTML}{8AC0D9}
\definecolor{PastelRed}{HTML}{00B0DA}
\definecolor{GoldenEarth}{HTML}{B4AD0C}
\colorlet{name}{black}
\colorlet{tagline}{PastelRed}
\colorlet{heading}{DarkPastelRed}
\colorlet{headingrule}{GoldenEarth}
\colorlet{subheading}{PastelRed}
\colorlet{accent}{PastelRed}
\colorlet{emphasis}{SlateGrey}
\colorlet{body}{LightGrey}

\renewcommand{\familydefault}{\sfdefault}

\name{\Large BOILEAU Guillaume}
\tagline{Ingénieur IVVQ | Intégration, Validation, Analyse technique et systèmes complexes}
\personalinfo{
  \email{guillaume.boileaupro@gmail.com}
  \phone{+33 6 59 94 85 84}
  \location{Alpes Maritimes, France}
  \linkedin{guillaume-boileau-phd-891b81265}
  \researchgate{Guillaume-Boileau-2}
  \orcid{0000-0002-3576-6968}
}

\setlength{\parskip}{0.75\baselineskip}

\begin{document}
\makecvheader
\vspace{-0.6cm}
\cvsection{Lettre de motivation}
\vspace{-0.4cm}

{\color{accent}\textbf{Objet :}} \textit{Candidature au poste d'Ingénieur IVVQ - F/H}

{\color{body}

Madame, Monsieur,

Je vous adresse ma candidature au poste d'Ingénieur IVVQ au sein de Scalian à Cannes. Ce poste correspond très directement à mon positionnement actuel : intervenir sur des systèmes complexes, contribuer à leur intégration, structurer les activités de vérification et de validation, analyser les écarts, assurer la traçabilité des constats et coordonner les échanges techniques nécessaires à la résolution des problèmes jusqu'à la conformité attendue.

Mon parcours m'a conduit à travailler sur des environnements particulièrement exigeants, dans lesquels la robustesse des résultats, la qualité de l'analyse, la compréhension des interactions entre sous-ensembles et la rigueur des méthodes de validation sont centrales. J'y ai développé une pratique solide de l'intégration système, de la validation, de l'analyse technique, du traitement d'incohérences, de la comparaison de comportements observés avec les attentes définies, ainsi que de la structuration de scénarios d'analyse et de campagnes de test.

J'ai notamment été amené à développer, structurer et valider des chaînes logicielles complexes, à comparer plusieurs pipelines ou sorties de traitements, à identifier les écarts significatifs, à documenter les anomalies et à proposer des cadres d'analyse permettant de rendre les résultats exploitables, traçables et défendables. Ce travail m'a demandé non seulement une forte rigueur technique, mais aussi une capacité à relier plusieurs niveaux de compréhension d'un système, depuis les hypothèses de départ et les exigences fonctionnelles jusqu'aux comportements finaux observés. C'est précisément cette capacité à faire le lien entre spécification, intégration, vérification et analyse de conformité que je souhaite aujourd'hui mettre au service d'un poste IVVQ.

Mon expérience récente chez VEV Platform Services France a renforcé cette approche dans un contexte plus directement industriel et opérationnel. J'y ai contribué au développement de services web et backend, à l'intégration de composants logiciels, au suivi de flux de données opérationnels et à l'analyse d'incidents impliquant des interactions entre logiciel et matériel. Cette expérience m'a permis de consolider ma capacité à intervenir sur des systèmes réels, à suivre le comportement global de services interconnectés, à investiguer des anomalies en production, à collaborer avec plusieurs intervenants techniques et à contribuer à la fiabilité des échanges entre briques système.

En parallèle, mes expériences précédentes m'ont permis de travailler sur des sujets nécessitant une forte maîtrise de la traçabilité, de la validation croisée, de la structuration des tests et de l'amélioration continue des méthodes. J'ai conçu et développé CATpyLISA from scratch comme cadre de comparaison, validation et consolidation de résultats complexes, avec un véritable enjeu de robustesse méthodologique, de cohérence globale et de qualité finale des produits obtenus. Ce type de travail m'a appris à ne pas seulement produire des résultats techniques, mais à construire le cadre qui permet de les tester, de les comparer, de les qualifier et de les rendre fiables dans la durée.

Je peux ainsi apporter à Scalian une contribution utile sur les dimensions clés du poste. Je suis à l'aise pour rédiger et dérouler des scénarios de test, participer à une stratégie de validation d'ensemble, vérifier la conformité d'un système vis-à-vis d'exigences définies, identifier et documenter les anomalies, suivre leur résolution avec les équipes concernées, contribuer à la définition de critères d'acceptation et maintenir le niveau de traçabilité nécessaire tout au long du processus d'intégration. J'ai également l'habitude de travailler sur des systèmes où plusieurs sous-ensembles interagissent, ce qui demande de conserver une vision d'ensemble, de coordonner les contributions et de partager un état de conformité lisible et exploitable.

Au-delà de la technique, je peux aussi apporter une réelle capacité de coordination. J'ai supervisé des stagiaires, animé des échanges scientifiques et techniques, coordonné des activités impliquant plusieurs contributeurs et structuré des démarches collectives autour d'objectifs complexes. J'ai appris à organiser le travail, à aligner des intervenants aux profils différents, à faire remonter les blocages de manière claire et à maintenir un cadre méthodique même sur des sujets difficiles. Cette capacité à porter un sujet collectivement me paraît particulièrement pertinente dans un environnement de conseil en ingénierie comme Scalian, où l'on attend à la fois rigueur, sens du service, adaptabilité et efficacité chez le client.

Ce qui m'intéresse particulièrement dans cette opportunité est justement la possibilité d'intervenir sur des projets techniquement riches, dans des environnements critiques et structurés, en mettant à profit un profil capable d'allier intégration, validation, analyse système, investigations techniques et coordination transverse. Mon parcours me permet d'être rapidement crédible sur des sujets complexes, d'apporter une méthode de travail rigoureuse, et de contribuer concrètement à la qualité, à la conformité et à la robustesse des systèmes sur lesquels j'interviens.

Rejoindre Scalian représenterait pour moi l'opportunité de mettre ces compétences au service de projets industriels exigeants, dans une logique d'IVVQ système, de service au client et de progression continue au sein d'une communauté d'experts. Je serais heureux de pouvoir échanger avec vous afin de vous présenter plus en détail ma motivation et la manière dont je pourrais contribuer efficacement à vos missions à Cannes.

Je vous prie d'agréer, Madame, Monsieur, l'expression de mes salutations distinguées.

\textbf{Guillaume Boileau}

}

\end{document}